\BukaLampiran

\lampiran{\textit{Source code}}

\lstset{
    basicstyle=\ttfamily\small,
    keywordstyle=\bfseries,
    commentstyle=\itshape\color{gray},
    breaklines=true,
    frame=single,
    numbers=left,
    numberstyle=\tiny,
    xleftmargin=2em,
    framexleftmargin=1.5em,
    tabsize=4,
}

\begin{lstlisting}[language=Python, caption={Algoritma \textit{Hybrid Vector-Graph Retrieval}}]
ALGORITHM: HybridVectorGraphRetrieval

INPUT:
    Q           : kueri pengguna
    VectorDB    : basis data vektor (Zilliz Cloud)
    GraphDB     : basis data graf (Neo4j AuraDB)
    Reranker    : model Cross-Encoder Reranker
    mode        : "naive" | "subgraph" | "hybrid"
    top_k       : jumlah dokumen yang diambil

OUTPUT:
    context     : konteks terstruktur untuk LLM

BEGIN
    result = {}

    # === TAHAP 1: Dekomposisi Kata Kunci ===
    IF mode IN {subgraph, hybrid} THEN
        keyword_clues = VectorDB.search("ContentKeyword", Q)
        (K_high, K_low) = LLM_decompose(Q, keyword_clues)
        local_query  = combine(K_low, Q)
        global_query = combine(K_high, Q)
    END IF

    # === TAHAP 2: Pencarian Vektor (Naive) ===
    IF mode IN {naive, hybrid} THEN
        paper_chunks = VectorDB.search("PaperChunk", Q, top_k * 3)
    END IF

    # === TAHAP 3: Pencarian Subgraf Lokal ===
    IF mode IN {subgraph, hybrid} THEN
        entities = EntitySubgraphRetrieval(local_query, VectorDB,
                                           GraphDB, top_k)
        subgraph_triples = GraphDB.shortest_path(entities, depth=2)
        text_units = VectorDB.search("PaperChunk", local_query)
    END IF

    # === TAHAP 4: Pencarian Relasi Global ===
    IF mode IN {hybrid} THEN
        relationships = VectorDB.search("RelationshipEmbedding",
                                         global_query, top_k * 2)
        FOR EACH rel IN relationships DO
            src_deg = GraphDB.get_degree(rel.src_id)
            tgt_deg = GraphDB.get_degree(rel.tgt_id)
            rel.distance = rel.distance /
                (1 + 0.05 * (log(src_deg+1) + log(tgt_deg+1)))
        END FOR
        relationships = sort_ascending(relationships, key=distance)
    END IF

    # === TAHAP 5: Fusi dan Pemeringkatan Ulang ===
    IF mode == hybrid THEN
        ranked_papers = RRF_Rerank(paper_chunks, entities,
                                    GraphDB, Reranker, Q, top_k)
    ELSE IF mode == naive THEN
        ranked_papers = Reranker.rerank(Q, paper_chunks, top_k)
    END IF

    # === TAHAP 6: Penyusunan Konteks ===
    context = format_context(
        paper_chunks  = ranked_papers,
        entities      = entities,
        relationships = relationships,
        subgraph      = subgraph_triples,
        text_units    = text_units
    )

    RETURN context
END
\end{lstlisting}
